\section{Method}
\label{sec:method}

% TODO(figure): add the pipeline figure \ref{fig:pipeline}:
%   (a) escalation attack schedule; (b) temporal edge encoder (3 branches + gate);
%   (c) GNN detection + edge-pruning remediation.

\subsection{Overview}
\label{sec:method_overview}
We study attacker detection and remediation in an LLM-based multi-agent system
(MAS) that engages in multiple rounds of dialogue. Building on the graph-based
detect-and-remediate paradigm~\citep{wang2025gsafeguard}, we make two
contributions that both hinge on the temporal structure of inter-agent
communication. Section~\ref{sec:method_prelim} formalizes the MAS, the
detection problem, and our threat model. Section~\ref{sec:method_attack}
introduces the \emph{escalation attack}, which exploits detectors that aggregate
an agent's utterances over time. Section~\ref{sec:method_tgat} presents
\textbf{T-GAT}, our temporal-aware detector whose \emph{temporal edge encoder}
preserves cross-turn dynamics, and Section~\ref{sec:method_train} describes
per-round detection, remediation, and training. Figure~\ref{fig:pipeline} gives
an overview.

\subsection{Preliminaries and Threat Model}
\label{sec:method_prelim}

\paragraph{Multi-agent system.}
We model a MAS as a directed graph $\mathcal{G}=(\mathcal{V},\mathcal{E})$ with
$N=|\mathcal{V}|$ agents. Given a user query, the agents communicate for $K$
rounds; at round $t$, agent $C_i$ produces an utterance $\mathbf{R}_i^{(t)}$
conditioned on the query and on the messages of its in-neighbors
$\mathcal{N}_{\text{in}}(C_i)$. A text encoder
$\mathcal{T}:\mathbb{T}\rightarrow\mathbb{R}^{D}$ (we use
MiniLM~\citep{wang2020minilm} via Sentence-Transformers~\citep{reimers2019sentencebert},
$D{=}384$) maps any utterance to an embedding.

\paragraph{Detection as node classification.}
Following~\citet{wang2025gsafeguard}, we cast attacker detection as node
classification on a \emph{multi-agent utterance graph}. A subset
$\mathcal{V}_{\text{atk}}\subseteq\mathcal{V}$ is compromised, and a detector
$\mathcal{D}$ must recover it from observed communication. Because a compromised
agent's toxic messages propagate, the compromised set grows across rounds,
$\mathcal{V}_{\text{atk}}^{(1)}\subseteq\cdots\subseteq\mathcal{V}_{\text{atk}}^{(K)}$,
so detection is run \emph{per round}. Once flagged, an agent is remediated by
pruning its outgoing edges so its messages no longer reach others.

\paragraph{Threat model.}
The attacker controls the utterances of the compromised agents
$\mathcal{V}_{\text{atk}}$ but cannot modify the detector, its training data, or
the honest agents. Its objective is twofold: \emph{(i)} maximize propagation of
malicious content through the MAS (high attack success rate), and \emph{(ii)}
minimize the probability of being flagged by $\mathcal{D}$. The attacker's only
lever is \emph{behavioral}: how it shapes its utterances across the $K$ rounds.
This is the regime where the temporal profile of an attack becomes a degree of
freedom that prior continuous attacks leave unused.

\subsection{Escalation Attack}
\label{sec:method_attack}

\paragraph{Motivation.}
Existing graph-based detectors summarize an agent's multi-turn behavior before
reasoning over the graph: the per-turn embeddings on an edge are collapsed into a
single vector by mean-pooling,
\begin{equation}
\label{eq:meanpool}
\bar{\mathbf{e}}_{ij}^{(t)}=\frac{1}{t}\sum_{\tau=1}^{t}\mathbf{z}_{ij}^{(\tau)},
\qquad \mathbf{z}_{ij}^{(\tau)}=\mathcal{T}\!\left(\mathbf{R}_{i\to j}^{(\tau)}\right),
\end{equation}
or by keeping only the last turn $\mathbf{z}_{ij}^{(t)}$. Either way, detection
depends on \emph{average} content and is invariant to \emph{how} content evolves.
This is the weakness our attack targets.

\paragraph{Design.}
The escalation attack assigns each compromised agent a monotonically increasing
\emph{intensity schedule} $\lambda_1\le\lambda_2\le\cdots\le\lambda_K$ with
$\lambda_1\approx 0$. At round $1$ the agent uses an instruction identical to that
of honest agents, making it indistinguishable at the start; over subsequent rounds
it gradually strengthens its push toward the malicious objective (light, then
medium, then full advocacy), while keeping surface form---tone, length, and
vocabulary---consistent to avoid abrupt per-turn jumps. Abstractly, the
attacker's per-round embedding can be written as a drift from the benign centroid
$\boldsymbol{\mu}_b$ toward a malicious centroid $\boldsymbol{\mu}_m$,
\begin{equation}
\label{eq:escalation}
\mathbf{z}^{(\tau)}_{\text{atk}}\approx(1-\lambda_\tau)\,\boldsymbol{\mu}_b
+\lambda_\tau\,\boldsymbol{\mu}_m .
\end{equation}

\paragraph{Why it works.}
Substituting Eq.~\eqref{eq:escalation} into the mean-pool of
Eq.~\eqref{eq:meanpool} gives
$\bar{\mathbf{e}}^{(t)}\approx\boldsymbol{\mu}_b+\bar{\lambda}^{(t)}
(\boldsymbol{\mu}_m-\boldsymbol{\mu}_b)$ with
$\bar{\lambda}^{(t)}=\tfrac{1}{t}\sum_{\tau\le t}\lambda_\tau$. Because the early
$\lambda_\tau$ are near zero, the running mean $\bar{\lambda}^{(t)}$ stays small
throughout the short horizon of a dialogue, so the aggregated feature remains
inside the benign region and below the detection threshold. A \emph{continuous}
attacker instead has $\lambda_\tau\approx1$ for all $\tau$, so its mean collapses
to $\boldsymbol{\mu}_m$ and is easily separated. The escalation attack thus
\emph{saturates} the aggregated feature with benign early turns, evading
aggregation-based detectors while its later, higher-intensity turns still
propagate malicious content to honest agents.

\subsection{T-GAT: Temporal-Aware Detection}
\label{sec:method_tgat}

\paragraph{Temporal edge features.}
T-GAT keeps the temporal axis that aggregation discards. The node feature is the
agent's initial-round embedding $\mathbf{x}_i=\mathcal{T}(\mathbf{R}_i^{(1)})$, and
each edge $(i,j)$ carries the full \emph{sequence} of per-turn embeddings
\begin{equation}
\label{eq:edgeseq}
\mathbf{Z}_{ij}^{(t)}=\big[\mathbf{z}_{ij}^{(1)},\dots,\mathbf{z}_{ij}^{(t)}\big]
\in\mathbb{R}^{t\times D},
\end{equation}
rather than a single pooled vector. The core of T-GAT is a \emph{temporal edge
encoder} $\mathcal{F}:\mathbb{R}^{t\times D}\rightarrow\mathbb{R}^{D}$ that maps
this sequence to a single edge embedding while retaining its dynamics.

\paragraph{Temporal edge encoder.}
Given a sequence $\mathbf{Z}\in\mathbb{R}^{t\times D}$ (we drop the $ij$ subscript),
we first add a learnable temporal positional encoding $\mathbf{P}$ and extract
three complementary views.
\emph{(1) Temporal self-attention} models which turns matter:
\begin{equation}
\mathbf{a}=\operatorname{mean}_\tau\,
\mathrm{LN}\!\big(\mathrm{MHA}(\tilde{\mathbf{Z}},\tilde{\mathbf{Z}},\tilde{\mathbf{Z}})
+\tilde{\mathbf{Z}}\big),\quad \tilde{\mathbf{Z}}=\mathbf{Z}+\mathbf{P}.
\end{equation}
\emph{(2) Temporal convolution} captures local sequential shifts between adjacent
turns, $\mathbf{c}=\mathrm{Conv1D}(\mathbf{Z})\in\mathbb{R}^{D}$.
\emph{(3) Cross-turn statistics} make escalation explicit:
\begin{equation}
\label{eq:stats}
\mathbf{s}=\mathrm{MLP}\big(\big[\,
\operatorname{mean}_\tau\mathbf{Z}\;\Vert\;\mathbf{z}^{(t)}\;\Vert\;
\operatorname{std}_\tau\mathbf{Z}\;\Vert\;
\textstyle\max_\tau\big|\mathbf{z}^{(\tau+1)}-\mathbf{z}^{(\tau)}\big|
\,\big]\big),
\end{equation}
where $\Vert$ denotes concatenation and the four terms are the temporal mean, the
last turn, the cross-turn standard deviation, and the largest step-wise change.
The last two terms are near zero for a consistent honest agent but large for an
escalating attacker, providing a direct signal that mean-pooling destroys. Finally, a learned gate
adaptively fuses content and temporal dynamics,
\begin{equation}
\label{eq:gate}
\mathbf{g}=\sigma\!\big(\mathbf{W}_g[\mathbf{a}\Vert\mathbf{c}\Vert\mathbf{s}]\big),
\quad
\mathbf{e}_{ij}=\mathbf{g}\odot\phi_c(\mathbf{a})
+(1-\mathbf{g})\odot\phi_t([\mathbf{c}\Vert\mathbf{s}]),
\end{equation}
where $\phi_c,\phi_t$ are linear projections and $\odot$ is the Hadamard product.
The gate lets the model rely on content for continuous attacks (where average
content is already discriminative) and on temporal dynamics for escalation
attacks (where it is not), so a single detector handles both regimes.

\subsection{Detection, Remediation, and Training}
\label{sec:method_train}

\paragraph{Graph detection and remediation.}
The encoded edge embeddings $\mathbf{e}_{ij}$ feed an $L$-layer edge-featured graph
attention network~\citep{velickovic2018gat} that propagates structural and
semantic dependencies. Writing $\mathcal{N}_i$ for the in-neighbors of $C_i$,
\begin{equation}
\mathbf{h}_i^{(l)}=\mathrm{COMB}\!\Big(\mathbf{h}_i^{(l-1)},\,
\mathrm{AGGR}\big\{\psi(\mathbf{h}_j^{(l-1)},\mathbf{e}_{ij}):
C_j\in\mathcal{N}_i\big\}\Big),
\end{equation}
after which each agent receives an attack probability
$p_i=\sigma\!\big(f_\theta(\mathbf{h}_i^{(L)})\big)$. At the end of round $t$, the
flagged set $\tilde{\mathcal{V}}_{\text{atk}}^{(t)}=\{C_i:p_i\ge 0.5\}$ is remediated
by pruning its outgoing edges,
\begin{equation}
\label{eq:prune}
\mathcal{E}^{(t+1)}\leftarrow\mathcal{E}^{(t+1)}\setminus
\bigcup_{C_i\in\tilde{\mathcal{V}}_{\text{atk}}^{(t)}}
\{e_{ij}\mid C_j\in\mathcal{V}\},
\end{equation}
so that flagged agents can no longer contaminate their neighbors in subsequent
rounds. Algorithm~\ref{alg:main} summarizes the defended dialogue loop.

\paragraph{Training objective.}
We train T-GAT with a binary cross-entropy loss over attack labels
$y_i\in\{0,1\}$,
\begin{equation}
\label{eq:loss}
\mathcal{L}=-\mathbb{E}_{C_i\sim\mathcal{V}}\big[
y_i\log p_i+(1-y_i)\log(1-p_i)\big].
\end{equation}
To make detection robust to the round at which it is invoked, we apply
\emph{temporal augmentation} during training: we randomly truncate each edge
sequence to its first $t'\le t$ turns, simulating early-stage detection when only
partial temporal evidence is available. This trains a single model to flag
attackers as early as possible rather than only after the full dialogue.

\begin{algorithm}[tb]
\caption{Defended multi-agent dialogue with T-GAT}
\label{alg:main}
\textbf{Input}: MAS $\mathcal{G}=(\mathcal{V},\mathcal{E})$, rounds $K$, trained
detector $\mathcal{D}_\theta$\\
\textbf{Output}: final answers and flagged set $\tilde{\mathcal{V}}_{\text{atk}}$
\begin{algorithmic}[1]
\STATE Run round $1$; encode utterances with $\mathcal{T}$.
\FOR{$t=1$ to $K$}
\STATE Build utterance graph with temporal edge features
$\mathbf{Z}_{ij}^{(t)}$ (Eq.~\ref{eq:edgeseq}).
\STATE $\mathbf{e}_{ij}\leftarrow\mathcal{F}(\mathbf{Z}_{ij}^{(t)})$ via the
temporal edge encoder (Eqs.~\ref{eq:stats}--\ref{eq:gate}).
\STATE Compute $p_i$ for all $C_i$; set
$\tilde{\mathcal{V}}_{\text{atk}}^{(t)}=\{C_i:p_i\ge0.5\}$.
\STATE Prune outgoing edges of $\tilde{\mathcal{V}}_{\text{atk}}^{(t)}$
(Eq.~\ref{eq:prune}).
\STATE Run round $t{+}1$ on the pruned topology.
\ENDFOR
\STATE \textbf{return} final answers, $\tilde{\mathcal{V}}_{\text{atk}}$
\end{algorithmic}
\end{algorithm}
